\par\Needspace{6\baselineskip}
\subsection{Hub ও Bridge-এর ঐতিহাসিক ধরন}
\textbf{পুরোনো হাবের ধরন:} Passive Hub মূলত সংযোগ ভাগ করে; সংকেত পুনর্গঠন করে না। Active Hub বিদ্যুৎ ব্যবহার করে সংকেত পুনর্গঠন করে। Intelligent Hub-এ ব্যবস্থাপনার সুবিধা থাকতে পারে। আধুনিক ল্যাবে সাধারণত সুইচ ব্যবহার করা হয়।
\textbf{ব্রিজের ব্যবহারভেদ:} Local Bridge কাছের অংশ এবং Remote Bridge দূরের অংশ যুক্ত করে। Wireless Bridge বেতার সংযোগ দিয়ে নেটওয়ার্কের অংশ সংযুক্ত করে। এই নামগুলো মাধ্যম ও অবস্থানের পরিচয় দেয়। MAC-এর ভিত্তিতে ব্রিজের মূল কাজ বদলায় না।

\subsection{Modem-এর অবস্থান ও Router-এর সহায়ক কাজ}
\begin{multicols}{2}
\textbf{Internal ও External Modem:} Internal Modem কম্পিউটারের ভেতরে কার্ড বা মডিউল হিসেবে থাকে। External Modem বাইরে থাকে এবং কেবল বা পোর্ট দিয়ে যুক্ত হয়। পুরোনো ডেস্কটপের ডায়াল-আপ মডেম কার্ড প্রথমটির এবং বাইরে রাখা USB মডেম দ্বিতীয়টির উদাহরণ। অবস্থান ও সংযোগের ভিত্তিতে এই শ্রেণি করা হয়। Modulation ও Demodulation-এর মূল কাজ একই থাকে।\par
\textbf{একটি Public IP-তে একাধিক সংযোগ:} NAT (Network Address Translation) IP Address রূপান্তর করে। PAT (Port Address Translation) বা NAT Overload-এ IP-এর পাশাপাশি পৃথক TCP/UDP Port Number দিয়ে এক Public IP-তে একাধিক কম্পিউটারের সংযোগ আলাদা রাখা হয়। ল্যাবের দুটি কম্পিউটার একই ওয়েবসাইট খুললেও রূপান্তরের সারণি অনুযায়ী উত্তর সঠিক কম্পিউটারে ফিরে যায়।\par
\end{multicols}
\textbf{DHCP (Dynamic Host Configuration Protocol):} DHCP সার্ভার স্বয়ংক্রিয়ভাবে IP Address ও প্রয়োজনীয় নেটওয়ার্ক কনফিগারেশন দেয়। বাসার রাউটারে এই সুবিধা থাকলে নতুন ফোনে ঠিকানা হাতে লিখতে হয় না। Routing, NAT/PAT ও DHCP আলাদা কাজ। সব রাউটারে সব সুবিধা থাকা আবশ্যক নয়। Default Gateway হিসেবে কাজ করলেই যন্ত্রটি প্রোটোকল রূপান্তরকারী হয়ে যায় না।
\textbf{DSL ও ADSL:} DSL = Digital Subscriber Line। টেলিফোনের তামার জোড়া তারে উপযুক্ত মডেম ও সংযোগযন্ত্র দিয়ে ডেটা সেবা দেওয়া যায়। কথার ও ডেটার সংকেত পৃথক ফ্রিকোয়েন্সি অঞ্চলে থাকলে ফিল্টার বা স্প্লিটার দিয়ে একই সংযোগে দুটিই ব্যবহার করা সম্ভব। ADSL (Asymmetric DSL)-এ ডাউনস্ট্রিম ও আপস্ট্রিমের বহনক্ষমতা অসমান; সাধারণত ডাউনলোডের ক্ষমতা বেশি। এটি সংযোগের প্রযুক্তি। Internal/External মডেমের অবস্থান এবং Star/Bus সংযোগের বিন্যাস বোঝায়।

\subsection{Modem, ONT ও Gateway নিয়ে বিভ্রান্তি}
\begin{tabularx}{\linewidth}{|p{.21\linewidth}|X|X|}
\hline
\textbf{নাম} & \textbf{মূল কাজ} & \textbf{একটির সঙ্গে অন্যটি গুলিয়ে ফেলবে না}\\\hline
Modem & ডেটা অনুযায়ী বাহক সংকেত বদলানো এবং প্রাপ্ত সংকেত থেকে ডেটা উদ্ধার করা; অর্থাৎ Modulation/Demodulation। & পুরোনো টেলিফোন মডেম দিয়ে ডিজিটাল–অ্যানালগ রূপান্তর বোঝানো যায়। আধুনিক সব ইন্টারনেট ইন্টারফেস একইভাবে রূপান্তর করে না।\\\hline
ONT & Optical Network Terminal: ফাইবার সংযোগের প্রান্তে আলোকসংকেত গ্রহণ করে স্থানীয় ডেটা ইন্টারফেস দেয়। & FTTH-তে ব্যবহৃত হয়। টেলিফোনের অ্যানালগ মডেম ও ONT একই যন্ত্র নয়।\\\hline
Default Gateway & স্থানীয় নেটওয়ার্কের বাইরের গন্তব্যে পাঠানোর পরবর্তী যন্ত্র; সাধারণত রাউটার। & কম্পিউটারের গেটওয়ে সেটিং মানেই প্রোটোকল রূপান্তরকারী নয়।\\\hline
Protocol Gateway & ভিন্ন প্রোটোকল বা অ্যাপ্লিকেশন ব্যবস্থার মধ্যে প্রয়োজনীয় রূপান্তর করে। & প্রোটোকল রূপান্তরের সাহায্যে ভিন্ন ব্যবস্থার মধ্যে যোগাযোগ সম্ভব হয়।\\\hline
\end{tabularx}
\textbf{একটি যন্ত্রে একাধিক কাজ:} বাসার তারবিহীন রাউটারে Routing, Ethernet Switch ও Wi-Fi Access Point-এর সুবিধা একত্রে থাকতে পারে। ফায়ারওয়ালও থাকতে পারে। তাই প্রতিটি কাজের জন্য পৃথক যন্ত্র আঁকা সবসময় প্রয়োজন নয়। তবে রাউটার নাম থেকেই সব মডেলে মডেম বা ONT আছে বলে ধরে নেওয়া যায় না।
\textbf{নিরাপত্তা:} ফায়ারওয়াল নির্ধারিত নীতি অনুযায়ী ট্রাফিক অনুমোদন বা বাধা দেয়। শক্তিশালী পাসওয়ার্ড, হালনাগাদ, ব্যাকআপ ও সচেতন ব্যবহারের প্রয়োজনও থাকে। সুইচ অপ্রয়োজনীয় ফ্রেম চলাচল কমায়। তবে শুধু সুইচ ব্যবহারের কারণে গোপন তথ্য স্বয়ংক্রিয়ভাবে এনক্রিপ্ট হয় না।


\textbf{যন্ত্রের পোর্ট ও দূরত্বের সংখ্যা:} ৮/১৬/২৪/৪৮-Port সুইচ পরিচিত কনফিগারেশন; এগুলো পোর্টসংখ্যার সংজ্ঞাগত সর্বোচ্চ সীমা নয়। ল্যাবের ১২ PC ও ১ Uplink-এর জন্য অন্তত ১৩টি উপযুক্ত পোর্ট প্রয়োজন। ডেটা হার ও অন্যান্য সুবিধা মিললে ১৬-Port মডেল নেওয়া যায়। রিপিটার বসানোর ব্যবধান কেবল বা বেতার মানের ওপর নির্ভর করে। ১৮৫ m পুরোনো 10BASE2 অংশের দৈর্ঘ্য; সব রিপিটারের জন্য একই ব্যবধান নয় (অংশ \ref{ch2:wired-numeric})।

\textbf{রাউটার ও গেটওয়ের ইতিহাস:} ১৯৭০-এর দশকে পৃথক নেটওয়ার্ক সংযুক্ত করার পরীক্ষায় Virginia Strazisar-এর প্রাথমিক TCP/IP Gateway/Router সফটওয়্যার গুরুত্বপূর্ণ অবদান রাখে। হাব, ব্রিজ, রিপিটার ও গেটওয়ে বিভিন্ন নকশায় বিকশিত হয়েছে। সব শ্রেণির জন্য একক উদ্ভাবক ও সাল নির্ধারণ নির্ভরযোগ্য নয়। Ethernet-এর উৎপত্তি নেটওয়ার্কের ইতিহাসে আছে। আধুনিক Ethernet সুইচ ও পুরোনো ভাগাভাগি করা বাস একই যন্ত্র নয়।
